\documentclass[10pt,letterpaper,sans]{moderncv}

\moderncvstyle{oldstyle}
\moderncvcolor{grey}

\usepackage[scale=0.88]{geometry}
\setlength{\hintscolumnwidth}{2.1cm}
\renewcommand*{\namefont}{\fontsize{24}{26}\bfseries\upshape}
\renewcommand*{\titlefont}{\Large\mdseries\upshape}
\renewcommand*{\addressfont}{\small\mdseries\upshape}

\name{Max}{Siegel}
\title{AI Research Scientist}

\email{maxs@mit.edu}

\begin{document}
\makecvtitle
\vspace{-1.1em}

\section{Profile}
\cvitem{}{AI research scientist at MIT specializing in computational cognitive science, multimodal reasoning, and perception. Builds and evaluates computational models of vision, sound, language, and social reasoning --- combining probabilistic methods with behavioral experiments to produce theory-driven insights relevant to AI.}

\section{Experience}
\cventry{2018--Present}{Research Scientist}{MIT, Department of Brain and Cognitive Sciences}{Cambridge, MA}{}{
\begin{itemize}
\item Build computational accounts of multimodal inference, physical reasoning, and perception using probabilistic modeling, generative modeling, and analysis-by-synthesis.
\item Design and run behavioral experiments to evaluate model predictions against human data, spanning 3D shape perception, auditory scene understanding, visual search, and object reasoning.
\item Take projects from initial question to publication: building stimuli, collecting data, running analysis, writing manuscripts, independently and in collaboration.
\item Work across cognitive science, computer science, neuroscience, and developmental psychology at MIT and with external collaborators.
\item Published 23 papers including 4 as co-first author; venues include \textit{Nature Human Behaviour}, \textit{Nature Communications}, and \textit{Journal of Experimental Psychology: General}.
\end{itemize}}

\section{Strengths}
\cvitem{Modeling}{Probabilistic modeling, generative modeling, analysis-by-synthesis, theory-driven model comparison}
\cvitem{Domains}{Multimodal reasoning, visual and auditory perception, concept learning, intuitive physics, human-AI interaction}
\cvitem{Methods}{Behavioral experiments, psychophysics, quantitative analysis, interdisciplinary collaboration}

\section{Selected Impact}
\cvitem{}{Co-first-authored \textit{Nature Human Behaviour} paper demonstrating that 3D shape perception is well-explained by a model integrating intuitive physics with analysis-by-synthesis.}
\cvitem{}{Co-first-authored \textit{Nature Communications} paper showing that children's exploratory play is precisely calibrated to the discriminability of competing hypotheses --- a quantitative link between cognition and behavior.}
\cvitem{}{Published additional work on auditory scene understanding, visual search, inverse rendering, face perception, and reasoning about hidden objects from multimodal cues.}

\section{Selected Publications}
\cvitem{2023}{Yildirim, I., \textbf{Siegel, M. H.}, Soltani, A. A., Ray Chaudhuri, S., \& Tenenbaum, J. B. Perception of 3D shape integrates intuitive physics and analysis-by-synthesis. \textit{Nature Human Behaviour}.}
\cvitem{2021}{\textbf{Siegel, M. H.}, Magid, R. W., Pelz, M., Tenenbaum, J. B., \& Schulz, L. E. Children's exploratory play tracks the discriminability of hypotheses. \textit{Nature Communications}.}

\section{Education}
\cventry{2018}{Ph.D. in Brain and Cognitive Sciences}{MIT}{Cambridge, MA}{}{}

\end{document}
